\documentclass[10pt,a4paper]{amsart}
\usepackage[margin=19mm]{geometry}
\usepackage{amsmath,amssymb,mathtools}
\usepackage[scr=rsfs,cal=euler]{mathalfa}
\usepackage{xfrac}
\usepackage[unicode,hidelinks,bookmarks=false]{hyperref}
\newcommand{\ie}{i.e.\ }
\newcommand{\cf}{cf.\ }
\newcommand{\resp}{respectively\ }
\newcommand{\Subsection}{\subsection}
\providecommand{\nobreakdash}{\nobreak\hbox{-}}
\setlength{\emergencystretch}{2em}
\multlinegap0pt
\usepackage{bm}
\usepackage{bbm}
\usepackage{dsfont}
\usepackage{xspace}
\usepackage{cancel}
\usepackage{mathtools}
\usepackage{amsthm}
\makeatletter
\@ifundefined{theorem}{\newtheorem{theorem}{Theorem}}{}
\@ifundefined{lemma}{\newtheorem{lemma}[theorem]{Lemma}}{}
\makeatother
\title[A 15/31 Counterexample Family to the Albertson-Berman Conjec]{A 15/31 Counterexample Family to the Albertson-Berman Conjecture}
\author{Heejae Jung}
\date{Source: arXiv:2608.17350v1 (CC BY 4.0)}
\begin{document}
\maketitle

\noindent\emph{Source and licence.} A 15/31 Counterexample Family to the Albertson-Berman Conjecture. Version arXiv:2608.17350v1 (CC BY 4.0), distributed under the Creative Commons Attribution 4.0 International licence, as recorded in the arXiv licence metadata for this exact version and shown on the abstract page https://arxiv.org/abs/2608.17350v1. Full rights notice: licenses/arxiv-2608.17350-CC-BY-4.0.txt.\par\smallskip

\noindent\textit{Editorial scope.} Counted unit: the Lemma whose source label is \texttt{lem:ico} in the article (source0.tex lines 239--247, source label \texttt{lem:ico}), delivered together with the article's own complete original proof of it (source0.tex lines 249--297, one \texttt{proof} environment of 49 lines). No mathematics is written, repaired or re-derived: statement and proof are the article's own text line for line. Required context retained uncounted: none. Every definition, display and label of those lines is retained unchanged. Omitted from this package and not redistributed: 1--238, 248--248, 298--1166. Nothing inside a retained excerpt is omitted or altered: every retained body is delivered exactly as the article's lines are written, so no omission receipt of any kind accompanies this package. Citations: the retained excerpts carry no citation command at all, so no bibliography entry of the article is transcribed and no reference is redistributed. Reference adaptation: none was needed and none was made; every reference inside the delivered unit resolves inside it, so no unresolved reference or citation is produced. Printed slips kept literal: no printed word, symbol, hypothesis or number of the counted statement or of its proof was altered. Layout and class: the article's own class is \texttt{article} with packages including geometry, amsmath, amssymb, amsthm, booktabs, array, tikz, hyperref; the wrapper is the lane's pinned \texttt{amsart} editorial wrapper at 10pt on A4, which restates the theorem-like environments the retained text needs and adds the article's own macro definitions that the retained text needs (none), with extra wrapper packages (bm, bbm, dsfont, xspace, cancel, mathtools). The delivered numbering is the unit's own. Assets: no retained span contains a figure environment, a graphics-inclusion command, a PDF-page inclusion, a table or a TikZ picture; no image file is copied into this package, the package declares no asset dependency and no image file is redistributed with it. Licence: version arXiv:2608.17350v1 of this article is distributed under the Creative Commons Attribution 4.0 International licence, as recorded in the arXiv licence metadata for this exact version and shown on the abstract page https://arxiv.org/abs/2608.17350v1. A full rights notice is delivered at licenses/arxiv-2608.17350-CC-BY-4.0.txt. Characters: every retained excerpt is pure ASCII, so the delivered engine, XeLaTeX with its default Latin Modern text font, reports no missing character.
\begin{lemma}[Icosahedral edge bounds]\label{lem:ico}
In the icosahedral graph $I$:
\begin{enumerate}
\item every five vertices span at least three edges;
\item every six vertices span at least five edges;
\item every six vertices containing the face $abf$ or the face $bfi$ span
      at least six edges.
\end{enumerate}
\end{lemma}

\begin{proof}
Fix a vertex $x$ of $I$.  Its six nonneighbours induce a wheel: one hub
joined to a $5$-cycle.  Each neighbour of $x$ has two neighbours on that
$5$-cycle.  These facts are the usual two-pentagon description of the
icosahedron and can also be read directly from the table above.  In
particular, the wheel has independence number two, so $\alpha(I)\le3$;
the independent set $\{a,i,k\}$ gives equality.

Let $R$ be a five-set.  If $I[R]$ had at most two edges, then, because
$\alpha(I)=3$, its two edges would have to be disjoint and the fifth vertex
$x$ would be isolated.  The other four vertices would lie in the
nonneighbour wheel of $x$ and would span only two edges.  This is impossible:
four rim vertices span three edges, while a set consisting of the hub and
three rim vertices spans at least three.  This proves (1).

Now let $R$ have six vertices.  If $e_I(R)\le4$, then $I[R]$ has a vertex
$x$ of degree at most one.  If its degree is zero, the other five vertices
lie in the nonneighbour wheel of $x$ and span at least five edges.  If its
degree is one, with neighbour $y$, the other four vertices span at least
three edges in that wheel.  Equality there forces four rim vertices.  But
$y$, being a neighbour of $x$, is adjacent to two rim vertices and hence to
at least one of the selected four.  Together with $xy$, this gives at least
five edges, a contradiction.  Thus (2) holds.

For (3), first use the face $F=\{a,b,f\}$.  The other nine vertices split
into
\[
 A=\{c,e,i\},\qquad B=\{d,j,m\},\qquad C=\{k,l,n\}.
\]
Each vertex of $A$ has two neighbours in $F$, each vertex of $B$ has one,
and vertices of $C$ have none.  The set $C$ induces a triangle, and the
$B$--$C$ edges are
\[
 dk,dl,jk,jn,ml,mn.
\]
Choose any three vertices outside $F$.  If $t$ of them lie in $C$, then:
for $t=0$ their attachments to $F$ contribute at least three edges; for
 $t=1$, the only tight attachment case chooses two vertices of $B$, and the
 chosen $C$-vertex meets at least one of them.  For $t=2$, the two chosen
 vertices of $C$ supply their mutual edge.  If the third vertex lies in $A$,
 its two attachments to $F$ complete the required three edges; if it lies in
 $B$, use its one attachment and at least one edge from it to the two chosen
 vertices of $C$.  For $t=3$, the triangle $C$ contributes three.  Adding the
 three edges of $F$ proves the claim for $abf$.  The automorphism
\[
 (a\ i)(b\ f)(c\ m)(d\ n)(e\ j)(k\ l)
\]
maps $abf$ to $bfi$, proving the other case.
\end{proof}

\end{document}
